\section{Introduction}
\label{sec:intro}

Multi-tenant LLM agents share runtimes, tool gateways, caches, and execution
capacity across tenants~\cite{mei2025aios,wu2023autogen}. This consolidation
makes permission checking a concurrency problem. A permission decision is
scheduled on an executor, reads policy state, may append an audit event, and
must complete while untrusted tenants generate work. Queueing, lock occupancy,
and executor rejection can therefore change the behavior of a permission
path even when resource admission and authorization are implemented as
separate modules. This interaction between concurrency and runtime control is
the focus of this paper.

We study one conditional but consequential failure mode. Suppose a runtime
sets a deadline for a permission decision and uses a permissive result when
that deadline expires. Hostile work that delays a shared policy executor can
then turn an executor-availability failure into an unsafe admission. The
condition is not inevitable: fail-closed handling preserves safety, and a
reserved permission executor can avoid the scheduling bottleneck. The point is
that deadline semantics, capacity allocation, queue pressure, and resource
signals must be evaluated together rather than hidden behind separate layer
interfaces.

\paragraph{Neuron.}
We present \textbf{Neuron}, a concurrency-aware runtime architecture with two core
mechanisms. First, permission checks execute on reserved capacity instead of
the executor used by tenant work. Second, resource admission exports a
pressure signal to the escalation policy; under sustained pressure, a
destructive request one level below the normal depth boundary is denied
instead of entering an interactive approval path. A per-turn trace identifier
joins the resource event and permission verdict for audit correlation.
Spawn-time identity propagation and privilege non-increase support the design,
while quotas and sandboxes are implementation mechanisms rather than
co-equal research contributions.

\paragraph{Systematic evaluation.}
An earlier draft used Python multiprocessing to model Java lock contention and
emphasized one 3\,ms deadline. That model demonstrated a possible mechanism
but could not establish behavior of the Neuron runtime. We remove it from the
evidence chain. The present study executes the Java~21 runtime
\texttt{RateLimitSyscallFilter} and \texttt{EscalationPolicy} classes in
Java~21, retains immutable JSONL observations, and treats each independently
started JVM as the replication unit. Five core runs are completed on each of
two environments, and a separate capacity experiment varies the permission
pool and queue. The four predeclared architectures are:

\begin{enumerate}
  \item a shared permission executor with fail-open timeout handling;
  \item the same shared executor with fail-closed handling;
  \item a reserved permission executor without cross-layer coupling; and
  \item reserved capacity plus Neuron's pressure-aware joint decision.
\end{enumerate}

The comparison separates two effects. If both reserved configurations meet
their deadlines, reservation explains availability. If only the coupled
configuration changes a depth-1 destructive request from \textsc{ask} to
\textsc{deny}, the difference is attributable to the resource signal rather
than executor isolation. A fixed-trace replay of previously collected API
sessions then tests the same decision path with application-level tool-call
sequences; it is a replay, not a new claim of live deployment.

\paragraph{Research questions.}
The paper follows the causal chain ``failure exists $\rightarrow$ mechanism
is isolated $\rightarrow$ behavior is checked outside one machine.''

\begin{description}
  \item[\textbf{RQ1}] Under what combinations of permission deadline,
  executor load, and lock-hold time does a shared Java policy executor miss
  its deadline, and how do timeout and capacity configurations differ?
  \item[\textbf{RQ2}] Does a resource-pressure signal change the escalation
  policy verdict beyond the effect of reserving executor capacity?
  \item[\textbf{RQ3}] Do the queueing and availability patterns reproduce on
  a second containerized environment, and how sensitive are they to
  permission-pool and queue capacity?
  \item[\textbf{RQ4}] Does the availability trade-off persist when the
  permission policy runs in a separate Java process behind an HTTP service
  with injected delay and pause faults?
  \item[\textbf{RQ5}] Does the concurrent decision path preserve completion
  for benign and adversarial API-derived tool-call traces, and what residual
  risks remain outside the local JVM/container boundary?
\end{description}

\paragraph{Contributions.}
This work contributes:

\begin{enumerate}
  \item a precise concurrency condition for executor-induced permission
  unavailability that states the required deadline and timeout semantics;
  \item a runtime architecture coupling permission-executor capacity with a
  resource-informed permission verdict; and
  \item a reproducible Java~21 evaluation with four fair baselines, two
  environments, independent JVM replications, deadline/load/capacity sweeps,
  automated raw-data aggregation, and a fixed API-trace replay.
\end{enumerate}

Identity propagation, privilege non-increase, and audit correlation are
supporting mechanisms. WASM, Docker, quota enforcement, and LLM routing are
engineering context and are reported in the artifact or appendix rather than
presented as independent innovations. AutoGen, LangGraph, and MCP are treated
as orchestration/protocol examples, not as competing multi-tenant security
systems; observations about their defaults therefore motivate integration
requirements but do not establish superiority over those projects.

Section~\ref{sec:bg} defines the failure conditions. Section~\ref{sec:threat}
states assumptions and non-goals. Section~\ref{sec:arch} describes the
architecture. Section~\ref{sec:eval} reports the cross-environment
evaluation, Section~\ref{sec:rw} positions the work, and Section~\ref{sec:conc}
concludes.
